\section{Introduction}\label{intro}

Building comfort diagnostics depend on how occupant state is represented. Predicted Mean Vote (PMV), expected thermal sensation, and binary comfort-band indicators are compact and useful, but each compresses a more heterogeneous response into a scalar or pass/fail state. That compression is consequential when the question is not only where a predicted response is centered, but how much probability is assigned to warm/hot or cool/cold categories and whether averaging across zones suppresses a localized signal.

Thermal sensation is not deterministic. Similar air temperature, radiant temperature, humidity, clothing, and metabolic conditions can produce different thermal sensation votes (TSVs) across people and contexts. Field studies attribute this dispersion to adaptation, expectation, physiology, recent exposure, behavior, and other unobserved differences \citep{deDear1998Adaptive,Brager2000Adaptive,Humphreys2002Adaptive,Foldvary2018GlobalDB}. A full class-probability distribution makes the modeled dispersion inspectable; a scalar mean does not show how its mass is allocated.

This paper treats the predicted seven-class TSV distribution as an information object. For classes $k\in\{-3,-2,-1,0,+1,+2,+3\}$, the predictor returns probabilities $p_k$. The expected sensation,
\begin{equation}
\mu_{\mathrm{TSV}}=\sum_{k=-3}^{3} k p_k,
\end{equation}
summarizes the center of that distribution. The TSV-tail probability,
\begin{equation}
p_{\mathrm{tail}}=\Pr(\mathrm{TSV}\le -2)+\Pr(\mathrm{TSV}\ge +2),
\end{equation}
summarizes the modeled probability mass in the cool/cold and warm/hot categories outside the central $-1,0,+1$ band. On the ASHRAE scale, $+1$ denotes ``slightly warm'' and $+2$ denotes ``warm''; the conventional PMV--PPD lineage associates the $\pm2$ and $\pm3$ sensation categories with dissatisfaction \citep{Fanger1970PMV,ASHRAE2023}. This convention gives the selected tail a discomfort-relevant interpretation, but it does not make sensation identical to satisfaction, preference, or acceptability. Accordingly, $p_{\mathrm{tail}}$ is a model-estimated TSV-category probability, not PPD or a measured dissatisfied fraction.

Two choices are kept separate throughout the analysis. The category boundary above defines the primary modeled event. The probability cutoff $p_{\mathrm{tail}}\ge0.20$ is only an illustrative reporting screen. It is not an ASHRAE or ISO criterion, a fitted optimum, a universal action threshold, or an 80\%-acceptability rule. The aggregation result is therefore examined over a range of probability screens, and an endpoint-only analysis using $p_{\mathrm{outermost}}=\Pr(\mathrm{TSV}\in\{-3,+3\})$ provides a stricter definition sensitivity.

Future weather is used as a controlled stress-test domain. The DOE Medium Office prototype, thermostat schedule, and present-day TSV-response model are held fixed while selected weather inputs change. The resulting experiment asks a conditional question: under the modeled building states generated by these weather files, what category-probability information is lost when the response distribution is replaced by a scalar, zones are averaged, or a single representative occupant input is assumed? It does not forecast how occupants, adaptive opportunities, building stock, or HVAC operation will evolve in the 2050s or 2080s.

The probabilistic and ordinal prediction framework is antecedent work rather than the novelty claimed here \citep{guo2026seven}. The contribution of this study is the diagnostic experiment:
\begin{itemize}
    \item a scalar-information audit comparing the full TSV distribution, its expected value, and PMV without treating PMV as a failed or obsolete comparator;
    \item a spatial-aggregation audit comparing mean-zone, both area-weighted operations, percentile-zone, and any-zone summaries;
    \item threshold, tail-definition, model-form, and weather-weighting robustness checks;
    \item a comfort-model input sensitivity showing what a single metabolic assumption suppresses; and
    \item a fixed-archetype stress test spanning 144 weather-selection roles, representing 119 unique source years across six cities, two emissions scenarios, and four time slices.
\end{itemize}

The paper does not validate an HVAC operating policy, optimize setpoints, estimate energy savings, or infer occupant exposure prevalence. Its claim is narrower: retaining calibrated TSV probabilities permits analysts to audit category mass, direction, spatial aggregation, and occupant-input sensitivity that scalar or building-average summaries do not directly expose.
